% Vaidyx benchmark paper. Build: docs/research/benchmark-paper/build.sh (tectonic). Tables are generated by build_data.py.
% Author names, e-mail addresses and the title are placeholders for the owner to edit.
\documentclass{article}
\usepackage{arxiv}
\usepackage[utf8]{inputenc}
\usepackage[T1]{fontenc}
\usepackage{lmodern}
\usepackage{microtype}
\usepackage{booktabs}
\usepackage{multirow}
\usepackage{makecell}
\usepackage{adjustbox}
\usepackage{graphicx}
\usepackage{array}
\usepackage{xcolor}
\usepackage{caption}
\usepackage{url}
\usepackage{hyperref}
\usepackage{natbib}

\captionsetup{font=small,labelfont=bf,skip=5pt}
\renewcommand{\theadfont}{\normalfont}
\renewcommand{\theadalign}{bc}
\newcommand{\fit}[1]{\adjustbox{max width=\textwidth}{#1}}
\hypersetup{colorlinks=true,linkcolor=black,citecolor=blue!55!black,urlcolor=blue!55!black}
\newcommand{\tabnote}[1]{\par\vspace{3pt}{\footnotesize #1\par}}
\newcommand{\pct}{\,\%}

\title{Full-Precision Benchmarking of Open Medical and General Language Models for a Tumour Board Assistant in India}

\renewcommand{\shorttitle}{Full-precision benchmarking of open medical and general models}
\renewcommand{\headeright}{Technical report}
\renewcommand{\undertitle}{Technical report}

\author{%
  Bharath Tangella \And Minmini Selvam \And Nikhileswara Rao Sulake
  \AND Phaneendra K Yalavarthy\thanks{Corresponding author.} \And Sai Manikanta Eswar Machara \And Usha Adiga
  \AND \normalfont TANUH, the AI Centre of Excellence in Healthcare, Indian Institute of Science, Bengaluru, India\\
  \texttt{\{bharath.tangella, minmini.selvam, nikhileswararao.sulake, phaneendra.yalavarthy,}\\
  \texttt{eswar.machara, usha.adiga\}@example.org}
}

\date{October 2026}

\hypersetup{
  pdftitle={Full-Precision Benchmarking of Open Medical and General Language Models for a Tumour Board Assistant in India},
  pdfauthor={Tangella, Selvam, Sulake, Yalavarthy, Machara, Adiga},
  pdfkeywords={medical language models, benchmark, tumour board, multimodal, grounding, India},
}

\begin{document}
\maketitle

\begin{abstract}
Vaidyx-TB is an assistant for hospital Tumour Boards in India. It checks, routes and summarises the records of a Case File and keeps every fact it shows tied to a verbatim quote from a source document; it does not interpret images, stage disease or advise treatment. Choosing the open-weight models behind it needs scores that are comparable with the published model cards and that reflect the exact weights the system runs. We benchmarked 13 open models (11 medical fine-tunes and 2 general models, 4B to 32B parameters, from developers in the USA, China and the UAE) on 11 public benchmarks, 5 text (9,995 questions) and 6 image (101,227 questions), with the published bf16 weights on one NVIDIA H200, greedy decoding, full test splits and no quantisation, under two protocols (step-by-step and direct answer): about one million graded answers. The harness reproduces published scores: MedGemma-27B-it lands within 2.5 points of its model card on all five text sets, and Lingshu-32B scores 82.4\pct{} against a published 83.4\pct{} on OmniMedVQA. Gemma-4-31B-it, a general model, is the strongest text model at 76.6\pct{} overall and 55.1\pct{} on MedXpertQA, where every medical model stays below 41\pct{}. HuatuoGPT-3-27B is the best medical model (71.2\pct{}) and leads OmniMedVQA (82.9\pct{}). Qwen3.5-4B matches MedGemma-27B-it on text at a seventh of the size, and Lingshu-32B leads radiology and pathology VQA. Full precision removes most of the gap seen in an earlier 4-bit run (older fine-tunes gain 11 to 30 points). We then report the grounded evaluations that decide what Vaidyx-TB ships: a reader and extractor bake-off on synthetic Case Files, 33 answerability probes with no fabrication from any of three candidate reasoners, 166 values extracted from 30 real TCGA pathology reports with none ungrounded, report-status detection measured on 9,523 real reports, a summary omission check, and an independent OCR number check (98.6\pct{} number recall; 99.4\pct{} of planted digit substitutions caught with no false alarm). Every answer is saved per question, so every number here can be re-scored.
\end{abstract}

\keywords{medical language models \and benchmark \and multimodal \and tumour board \and grounding \and India}

\section{Introduction}

A Tumour Board (multidisciplinary tumour board, MDT) decides the treatment plan of a cancer patient from the records of several departments. In India most of that preparation is manual. A 2026 survey of Indian cancer centres reports that 88.3\pct{} keep Board records on paper or in mixed form, that 53.3\pct{} have no radiologist taking part in the Board, and that 48.2\pct{} have no system to monitor whether the recommended treatment was given \citep{sarveswaran2026}.\footnote{Two of the three figures are derived from the survey rather than stated in it: 88.3\pct{} is 63.5\pct{} physical plus 24.8\pct{} mixed records, and 53.3\pct{} is the complement of the 46.7\pct{} radiologist participation it reports.} Vaidyx-TB, built at TANUH, addresses the preparation: it reads the documents of a Case File (pathology, radiology, laboratory and discharge reports, scanned or typed), checks the hospital's pre-authorisation requirements, extracts the values the Board needs, routes requests to Pathology and Radiology, drafts the case summary and the specialty briefs, and records the decision and the follow-up. Its scope is fixed by design: it checks, routes, summarises and flags. It never interprets an image, computes a stage or advises a treatment, and every value or sentence it shows as fact carries a verbatim quote from the source document; anything that cannot be verified is dropped or marked for review.

The system runs open-weight models on a hospital's own server. Two questions decide which models: how much medical knowledge a model has, measured on public examinations, and how reliably it stays grounded in the documents it is given, which no public benchmark measures. This report answers the first question with a benchmark that we could trust to be like for like with the published model cards, and the second with Vaidyx's own evaluations.

The benchmark differs from the usual model-card table in four ways. (1) Every model runs with its exact published bf16 weights; no quantisation or other lossy step sits anywhere in the pipeline, because the deployed system also runs bf16. (2) Every dataset is run in full; nothing is sub-sampled. (3) Every model runs under the same two protocols, a step-by-step protocol that matches how model cards report their scores and a direct-answer protocol that measures knowledge without reasoning and the discipline to follow a strict output format. (4) Every answer is saved per question, so every score can be re-checked or re-scored from the saved replies without regenerating anything. An earlier internal run of seven of these models used 4-bit weights (Q4\_K\_M) on a small workstation; this report shows how much of that run's gap was the compression.

Our contributions are:
\begin{itemize}
  \item a like-for-like benchmark of 13 open medical and general models on 11 public text and image benchmarks (111,222 questions, about one million graded answers), with the protocol, the scoring rules and the comparison with every published number we could find (Sections~\ref{sec:models} to \ref{sec:throughput});
  \item evidence that the harness reproduces published scores where the protocol matches, and a precise account of where it does not and why (Tables~\ref{tab:text-pub} and \ref{tab:img-pub});
  \item a measurement of what full precision changes against a 4-bit run of the same weights (Table~\ref{tab:q4});
  \item the grounded evaluations that select Vaidyx-TB's reader, extractor and reasoner, measured on synthetic Case Files and on real pathology reports (Section~\ref{sec:grounded}).
\end{itemize}

\section{Models}
\label{sec:models}

Table~\ref{tab:models} lists the 13 models. Each row carries two labels that appear in every result table of this report and in the Vaidyx-TB user interface: the developer's country, and whether the model is a medical fine-tune or a general model that is not medically validated. Seven models read images; the other six are text only and appear in the text tables only. Seven were part of the earlier 4-bit report. All weights are the published bf16 checkpoints from Hugging Face. Three models (HuatuoGPT-3, Baichuan-M2 and Qwen3.5) have a thinking mode; it was switched off, so that every model sees the same visible step-by-step prompt and the same token budget. Two models of the earlier report, a medical mixture-of-experts model and Asclepius-8B, were not repeated because no bf16 checkpoint was in this queue.

\begin{table}[htbp]
\centering\footnotesize\setlength{\tabcolsep}{5pt}
\caption{The 13 models. Type: Medical = fine-tuned for medicine; Non-medical = a general model, not medically validated. Earlier report = part of the 4-bit run of September 2026.}
\label{tab:models}
\fit{\begin{tabular}{@{}llllllc@{}}
\toprule
Model & Params & Developer & Country & Type & Reads images & Earlier report \\
\midrule
Gemma-4-31B-it & 31B & Google DeepMind & USA & Non-medical & Yes & New \\
HuatuoGPT-3-27B & 27B & FreedomIntelligence, CUHK-Shenzhen & China & Medical & Yes & New \\
Qwen3.5-4B & 4B & Alibaba Qwen team & China & Non-medical & Yes & New \\
MedGemma-27B-text-it & 27B & Google & USA & Medical & No & New \\
MedGemma-27B-it & 27B & Google & USA & Medical & Yes & Yes \\
Baichuan-M2-32B & 32B & Baichuan AI & China & Medical & No & New \\
Lingshu-32B & 32B & Alibaba DAMO Academy & China & Medical & Yes & New \\
MedGemma-4B-it & 4B & Google & USA & Medical & Yes & Yes \\
Med42-8B (v2) & 8B & M42 Health & UAE & Medical & No & Yes \\
JSL-MedLlama-3-8B-v2.0 & 8B & John Snow Labs & USA & Medical & No & Yes \\
MedGemma-1.5-4B-it & 4B & Google & USA & Medical & Yes & Yes \\
OpenBioLLM-8B & 8B & Saama AI Labs & USA & Medical & No & Yes \\
MedLlama2-7B & 7B & Community fine-tune of Llama 2 & USA & Medical & No & Yes \\
\bottomrule
\end{tabular}
}
\end{table}

\section{Benchmarks}
\label{sec:data}

Table~\ref{tab:datasets} gives a data card for each of the 11 benchmarks: five text sets (9,995 questions) and six image sets (101,227 questions). All are run on their full test split; nothing is sub-sampled. MedQA, MedMCQA, PubMedQA and MMLU-Med are the four text sets used on the MedGemma, Med42 and OpenBioLLM model cards, which makes the comparison in Table~\ref{tab:text-pub} possible. MedMCQA, drawn from the Indian postgraduate entrance examinations, is the most India-relevant text set and the largest, so it carries the most weight in the overall text score, which is the accuracy over all 9,995 questions together. MedXpertQA is filtered so that strong models fail it; its ten options put chance at 10\pct{}. The six image sets mix closed questions (yes/no, 2 to 5 options) and open short answers; OmniMedVQA alone has 88,996 questions over 42 public datasets and eight modalities.

\begin{table}[htbp]
\centering\footnotesize\setlength{\tabcolsep}{4pt}
\caption{Data cards of the 11 benchmarks. Text sets above the rule (9,995 questions), image sets below (101,227 questions). Chance = the score of random guessing.}
\label{tab:datasets}
\begin{tabular}{@{}>{\raggedright\arraybackslash}p{2.2cm}>{\raggedright\arraybackslash}p{6.6cm}>{\raggedright\arraybackslash}p{3.0cm}>{\raggedright\arraybackslash}p{2.2cm}>{\raggedright\arraybackslash}p{1.3cm}@{}}
\toprule
Dataset & What it tests & Questions & Format & Chance \\
\midrule
MedQA & US medical licensing exam (USMLE) style clinical vignettes: diagnosis, next step, mechanism & 1,273 (test) & 4 options & 25 \% \\
MedMCQA & Indian postgraduate medical entrance exams (AIIMS, NEET-PG) across 21 subjects; the most India-relevant text test & 4,183 (validation; test answers are hidden) & 4 options & 25 \% \\
PubMedQA & Reading a research abstract and answering a yes / no / maybe question about it (evidence reading, not recall) & 1,000 (expert-labelled) & yes / no / maybe & 33 \% \\
MMLU-Med & Six medical subjects of MMLU: anatomy, clinical knowledge, college biology, college medicine, medical genetics, professional medicine & 1,089 (test) & 4 options & 25 \% \\
MedXpertQA & Expert-level, deliberately hard questions from 17 specialty board exams, filtered so that strong models fail them; text subset & 2,450 (test) & 10 options & 10 \% \\
\midrule
VQA-RAD & Radiology images (CT, MRI, X-ray) with clinician-written questions & 451 (test) & yes/no and short answer & varies \\
SLAKE & Radiology images with organ, position and abnormality questions (English subset) & 1,061 (test) & yes/no and short answer & varies \\
PathVQA & Pathology images and textbook figures; open answers are long medical phrases, so exact match is strict & 6,719 (test) & yes/no and short answer & varies \\
PMC-VQA & Figures from PubMed Central papers (all modalities), cleaned test set & 2,000 (test\_clean) & 4 options & 25 \% \\
MedXpertQA-MM & Expert board-exam questions that need one or more clinical images plus a vignette & 2,000 (test) & 5 options & 20 \% \\
OmniMedVQA & Open-access part of OmniMedVQA: 42 public datasets, 8 modalities (MRI, CT, ultrasound, X-ray, dermoscopy, microscopy, fundus, OCT); mostly modality, anatomy and disease recognition & 88,996 & 2 to 4 options & about 27 \% \\
\bottomrule
\end{tabular}

\end{table}

\section{Protocol}
\label{sec:protocol}

\paragraph{Hardware and engine.} One NVIDIA H200 NVL (141\,GB), vLLM 0.30 \citep{kwon2023vllm}, PyTorch 2.13, CUDA 13. One model on the GPU at a time, 90\pct{} of memory, prefix caching on. The GPU ran power-capped at 300\,W after earlier hardware faults; the cap slows runs but changes no answer, since decoding is greedy.

\paragraph{Weights and decoding.} Exact published bf16 weights; no quantisation or other lossy compression anywhere. Greedy decoding (temperature 0), seed 0. Each model uses its own chat template; thinking mode off where a switch exists.

\paragraph{Two protocols.} Under the \emph{step-by-step} protocol (chain of thought, CoT) the instruction is: ``Answer the following multiple-choice question. Think through the question step by step, then finish with a final line of the form `Answer: X' where X is the letter of the correct option.'' The model may generate up to 2,048 new tokens, and the answer is read from the final \texttt{Answer:} line. This is how model cards report their results, so these are the scores to compare with published numbers. Under the \emph{direct} protocol the instruction is: ``Answer the following multiple-choice question. Respond with only the letter of the correct option.'', the prompt ends with \texttt{Answer:}, and the model may generate at most 32 new tokens (yes/no, or a short phrase for open image questions). It measures knowledge without reasoning and, at the same time, whether the model follows a strict output format. The question and its options follow the instruction in the same user turn, through each model's own chat template.

\paragraph{Scores.} The \emph{strict} score counts an answer only if a valid option can be read from the reply; a model that starts explaining and is cut off at 32 tokens scores zero for that question even if it was heading to the right answer. The \emph{first-token} score is a format-free check on the direct protocol: the option letter with the highest probability among the top-20 candidates for the first generated token. It shows what the model knew when the strict score is dragged down by formatting. \emph{Parsed\pct{}} is the share of replies from which an answer could be read; below about 95\pct{} a score is partly a formatting result. For VQA-RAD, SLAKE and PathVQA the score is split into closed (yes/no) and open short-answer questions; open answers are scored by exact match after normalising case and punctuation, which is strict and penalises every model equally. \emph{Errors} are questions the serving engine refused (an image too large for the model's context window, or more than four images); they count as wrong.

\paragraph{Published numbers.} Every published score was opened directly from the official model card or paper (Section~\ref{sec:sources}, pages opened on 9 October 2026). Where a card uses a different scoring rule, the rule is stated next to the number.

\paragraph{The earlier run.} The previous internal report (September 2026) ran seven of these models as Q4\_K\_M 4-bit weights with Ollama on an NVIDIA GB10 workstation under a direct letter prompt, on the same five text splits. Its direct prompt is closest to this run's direct protocol.

\paragraph{Scale and safety of the run.} 13 models under two protocols on text, 7 vision models on 6 image sets, plus the step-by-step protocol on the two multiple-choice image sets: about 1.0 million graded answers. OmniMedVQA alone took about 14 hours across the seven vision models. A thermal guard and GPU telemetry ran throughout. Every answer is saved per question (JSONL), so runs resume after an interruption and any score can be re-scored from the saved replies.

\section{Text results}
\label{sec:text}

\subsection{Step-by-step accuracy}

Table~\ref{tab:text-cot} ranks the 13 models by their overall step-by-step accuracy. Gemma-4-31B-it, a general model, leads every column except PubMedQA: MedQA 93.2, MedMCQA 78.5, MMLU-Med 94.6, and 55.1 on MedXpertQA, where every medical model stays below 41. HuatuoGPT-3-27B is the best medical model at 71.2\pct{} overall and has the highest PubMedQA score (81.4). Qwen3.5-4B, a 4B general model, reaches 62.4\pct{}, ahead of MedGemma-27B-it (61.7\pct{}) and within one point on MedQA (82.2 against 83.1). MedMCQA is the hardest of the four standard sets: the seven strongest models score 9 to 16 points lower on it than on MedQA. Two scores are held down by format rather than knowledge, which the last column shows: Med42-8B parses only 75\pct{} of its PubMedQA replies and MedGemma-1.5-4B only 71\pct{} of its MedQA replies.

\begin{table}[htbp]
\centering\small\setlength{\tabcolsep}{4pt}
\caption{Step-by-step (CoT) accuracy in \pct{}, ranked by overall accuracy over all 9,995 text questions. Bold = best in column. Min parsed = the lowest parse rate across the five sets.}
\label{tab:text-cot}
\fit{\begin{tabular}{@{}rlllrrrrrrr@{}}
\toprule
\# & Model & Origin & Type & MedQA & MedMCQA & PubMedQA & MMLU-Med & MedXpertQA & Overall & Min parsed \% \\
\midrule
1 & Gemma-4-31B-it & USA & Non-medical & \textbf{93.2} & \textbf{78.5} & 80.7 & \textbf{94.6} & \textbf{55.1} & 76.6 & 100.0 \\
2 & HuatuoGPT-3-27B & China & Medical & 91.0 & 75.3 & \textbf{81.4} & 92.5 & 40.4 & 71.2 & 97.7 \\
3 & Qwen3.5-4B & China & Non-medical & 82.2 & 68.0 & 78.7 & 88.4 & 24.3 & 62.4 & 96.9 \\
4 & MedGemma-27B-text-it & USA & Medical & 80.2 & 71.5 & 74.7 & 88.9 & 19.3 & 62.1 & 81.8 \\
5 & MedGemma-27B-it & USA & Medical & 83.1 & 68.5 & 76.9 & 87.6 & 21.2 & 61.7 & 92.4 \\
6 & Baichuan-M2-32B & China & Medical & 81.3 & 69.4 & 68.7 & 88.9 & 21.4 & 61.2 & 99.9 \\
7 & Lingshu-32B & China & Medical & 76.7 & 62.7 & 80.3 & 84.0 & 20.2 & 58.2 & 99.5 \\
8 & MedGemma-4B-it & USA & Medical & 66.1 & 56.1 & 75.8 & 71.3 & 14.4 & 50.8 & 99.0 \\
9 & Med42-8B (v2) & UAE & Medical & 62.7 & 58.6 & 57.9 & 73.7 & 13.8 & 49.7 & 74.9 \\
10 & JSL-MedLlama-3-8B-v2.0 & USA & Medical & 58.8 & 55.5 & 76.3 & 72.9 & 12.5 & 49.4 & 98.2 \\
11 & MedGemma-1.5-4B-it & USA & Medical & 46.9 & 50.5 & 60.9 & 64.7 & 10.2 & 42.7 & 70.6 \\
12 & OpenBioLLM-8B & USA & Medical & 49.9 & 48.0 & 52.6 & 67.5 & 11.8 & 41.9 & 91.4 \\
13 & MedLlama2-7B & USA & Medical & 38.2 & 34.9 & 69.9 & 48.3 & 10.5 & 34.3 & 95.8 \\
\bottomrule
\end{tabular}
}
\end{table}

\subsection{Direct answers and the gain from reasoning}

Table~\ref{tab:text-direct} gives the direct-answer scores, the first-token score, the parse rate and the gain from reasoning (step-by-step minus direct, in points). HuatuoGPT-3-27B is the best direct-answer model (66.0\pct{}), which matters for a system that asks short, structured questions. Reasoning helps the general models most: Gemma-4-31B-it gains 17.4 points overall and 44.1 points on MedXpertQA (11.1 to 55.1), and Qwen3.5-4B gains 9.3 points. Two Google models cannot follow the letter-only instruction: MedGemma-27B-text-it scores 8.0 on MedQA under the direct protocol with 25\pct{} of replies parseable, against 80.2 with reasoning, and MedGemma-1.5-4B behaves the same way; their first-token scores (46.7 and 35.7 overall) show what they knew. Baichuan-M2's direct replies all began with a stray \texttt{</think>} tag that the parser rejected; the saved replies were re-scored with the tag removed, and nothing was regenerated (dagger in the table).

\begin{table}[htbp]
\centering\footnotesize\setlength{\tabcolsep}{3.5pt}
\caption{Direct-answer accuracy in \pct{} (at most 32 new tokens), ranked by overall. First-token = the format-free score. Parsed = the lowest parse rate across the five sets. CoT gain = overall step-by-step minus overall direct, in points. $^{\dagger}$Re-scored from the saved replies after removing a stray thinking tag.}
\label{tab:text-direct}
\fit{\begin{tabular}{@{}lllrrrrrrrrr@{}}
\toprule
Model & Origin & Type & MedQA & MedMCQA & PubMedQA & MMLU-Med & MedXpertQA & Overall & \thead{First-token\\overall} & \thead{Parsed\\\%} & \thead{CoT\\gain} \\
\midrule
HuatuoGPT-3-27B & China & Medical & 83.7 & 70.4 & 80.8 & 90.2 & 32.6 & 66.0 & 66.0 & 99.9 & +5.2 \\
Lingshu-32B & China & Medical & 74.3 & 65.8 & 79.2 & 86.0 & 22.7 & 59.8 & 59.8 & 100.0 & -1.7 \\
Gemma-4-31B-it & USA & Non-medical & 68.4 & 71.7 & 79.6 & 90.2 & 11.1 & 59.2 & 68.8 & 21.8 & +17.4 \\
Baichuan-M2-32B$^{\dagger}$ & China & Medical & 74.6 & 67.1 & 71.7 & 86.9 & 17.9 & 58.6 & -- & -- & +2.6 \\
MedGemma-27B-it & USA & Medical & 66.3 & 60.7 & 77.2 & 81.0 & 12.3 & 53.4 & 55.5 & 70.4 & +8.3 \\
Qwen3.5-4B & China & Non-medical & 67.6 & 58.0 & 76.5 & 79.0 & 16.2 & 53.1 & 53.1 & 100.0 & +9.3 \\
Med42-8B (v2) & UAE & Medical & 62.1 & 60.1 & 77.2 & 72.9 & 15.1 & 52.4 & 52.5 & 99.9 & -2.7 \\
JSL-MedLlama-3-8B-v2.0 & USA & Medical & 45.0 & 56.5 & 77.2 & 67.9 & 10.0 & 47.0 & 47.8 & 74.5 & +2.4 \\
MedGemma-4B-it & USA & Medical & 55.2 & 52.0 & 74.9 & 64.9 & 13.2 & 46.6 & 46.6 & 100.0 & +4.2 \\
OpenBioLLM-8B & USA & Medical & 43.7 & 40.3 & 44.4 & 58.2 & 9.1 & 35.4 & 39.2 & 59.9 & +6.5 \\
MedGemma-27B-text-it & USA & Medical & 8.0 & 38.2 & 72.3 & 49.0 & 2.0 & 30.1 & 46.7 & 14.2 & +32.0 \\
MedGemma-1.5-4B-it & USA & Medical & 17.8 & 40.7 & 56.1 & 41.5 & 1.7 & 29.8 & 35.7 & 16.3 & +12.9 \\
MedLlama2-7B & USA & Medical & 16.4 & 33.2 & 68.0 & 40.3 & 4.1 & 28.2 & 33.3 & 38.6 & +6.1 \\
\bottomrule
\end{tabular}
}
\end{table}

\subsection{Comparison with published scores}

Table~\ref{tab:text-pub} sets our step-by-step scores against the published ones for the eight models that report any of these five sets. MedGemma-27B-it lands within 2.5 points of its card on all five sets (mean $-1.1$), MedGemma-4B-it slightly above its card (mean $+1.2$), and Lingshu-32B within 3.4 points on every set (mean $-0.4$); these three confirm that the harness is like for like with a generated-answer protocol. MedGemma-27B-text-it is 3.4 points below its card on average, mostly on MedQA and MedXpertQA. MedGemma-1.5-4B is 10 points below its card, mainly because 29\pct{} of its MedQA replies never reach a final \texttt{Answer:} line within 2,048 tokens. Med42, JSL-MedLlama and OpenBioLLM are scored on their cards by answer likelihood with the lm-evaluation-harness \citep{gao2023harness}, not by generated answers, which usually gives a few points more than a generated-answer test; OpenBioLLM is 12.3 points below its card on average and its PubMedQA gap is 21.5 points. No published scores on these sets exist for Gemma-4-31B, Qwen3.5-4B, HuatuoGPT-3-27B, Baichuan-M2-32B or MedLlama2-7B.

\begin{table}[htbp]
\centering\footnotesize\setlength{\tabcolsep}{4pt}
\caption{This run (bf16, step-by-step) against the published scores, in \pct{}. Difference = this run minus published, in points; the mean is over the sets the source reports.}
\label{tab:text-pub}
\fit{\begin{tabular}{@{}llrrrrrr@{}}
\toprule
Model & Row & MedQA & MedMCQA & PubMedQA & MMLU-Med & MedXpertQA & \thead{Mean\\diff.} \\
\midrule
\multirow{3}{*}{MedGemma-27B-text-it} & This run (bf16, CoT) & 80.2 & 71.5 & 74.7 & 88.9 & 19.3 & \\
 & Published (Google card) & 87.7 & 74.2 & 76.8 & 87.0 & 25.7 & \\
 & Difference & -7.5 & -2.7 & -2.1 & +1.9 & -6.4 & -3.4 \\
\midrule
\multirow{3}{*}{MedGemma-27B-it} & This run (bf16, CoT) & 83.1 & 68.5 & 76.9 & 87.6 & 21.2 & \\
 & Published (Google card) & 85.3 & 70.2 & 77.2 & 86.2 & 23.7 & \\
 & Difference & -2.2 & -1.7 & -0.3 & +1.4 & -2.5 & -1.1 \\
\midrule
\multirow{3}{*}{Lingshu-32B} & This run (bf16, CoT) & 76.7 & 62.7 & 80.3 & 84.0 & 20.2 & \\
 & Published (Lingshu report) & 74.7 & 66.1 & 77.8 & 84.7 & 22.7 & \\
 & Difference & +2.0 & -3.4 & +2.5 & -0.7 & -2.5 & -0.4 \\
\midrule
\multirow{3}{*}{MedGemma-4B-it} & This run (bf16, CoT) & 66.1 & 56.1 & 75.8 & 71.3 & 14.4 & \\
 & Published (Google card) & 64.4 & 55.7 & 73.4 & 70.0 & 14.2 & \\
 & Difference & +1.7 & +0.4 & +2.4 & +1.3 & +0.2 & +1.2 \\
\midrule
\multirow{3}{*}{Med42-8B (v2)} & This run (bf16, CoT) & 62.7 & 58.6 & 57.9 & 73.7 & 13.8 & \\
 & Published (M42 card$^{a}$) & 62.8 & 61.3 & -- & 75.8 & -- & \\
 & Difference & -0.1 & -2.7 & -- & -2.1 & -- & -1.6 \\
\midrule
\multirow{3}{*}{JSL-MedLlama-3-8B-v2.0} & This run (bf16, CoT) & 58.8 & 55.5 & 76.3 & 72.9 & 12.5 & \\
 & Published (JSL card$^{b}$) & 61.4 & 61.2 & 74.2 & 77.6 & -- & \\
 & Difference & -2.6 & -5.7 & +2.1 & -4.7 & -- & -2.7 \\
\midrule
\multirow{3}{*}{MedGemma-1.5-4B-it} & This run (bf16, CoT) & 46.9 & 50.5 & 60.9 & 64.7 & 10.2 & \\
 & Published (Google card) & 69.1 & 59.8 & 68.2 & 69.6 & 16.4 & \\
 & Difference & -22.2 & -9.3 & -7.3 & -4.9 & -6.2 & -10.0 \\
\midrule
\multirow{3}{*}{OpenBioLLM-8B} & This run (bf16, CoT) & 49.9 & 48.0 & 52.6 & 67.5 & 11.8 & \\
 & Published (Saama card$^{b}$) & 59.0 & 56.9 & 74.1 & 77.1 & -- & \\
 & Difference & -9.1 & -8.9 & -21.5 & -9.6 & -- & -12.3 \\
\bottomrule
\end{tabular}
}
\tabnote{The Med42, JSL and Saama cards score by answer likelihood (lm-evaluation-harness), not by generated answers. $^{a}$MMLU covers the clinical topics. $^{b}$MMLU is the mean of the six medical subjects.}
\end{table}

\subsection{Full precision against the earlier 4-bit run}

Table~\ref{tab:q4} compares the seven models of the September 2026 report, run then as Q4\_K\_M weights on a GB10 workstation, with this run on the same five splits. With the same weights at full precision and allowed to reason, MedGemma-27B-it rises from 69.8 to 83.1 on MedQA, and the older 7B to 8B fine-tunes gain between 10.6 and 30.1 points on average over the five sets (JSL-MedLlama $+30.1$, OpenBioLLM $+20.6$, MedGemma-1.5-4B $+15.9$, MedLlama2 $+10.6$). Under the direct protocol, which is closest to the earlier prompt, the gains are smaller and in two cases negative, so most of the earlier gap was the combination of 4-bit compression and a short-answer format, not the models themselves.

\begin{table}[htbp]
\centering\footnotesize\setlength{\tabcolsep}{4pt}
\caption{The earlier 4-bit run against this run, in \pct{}. Mean gain = the average of the five scores minus the Q4 average, in points.}
\label{tab:q4}
\fit{\begin{tabular}{@{}llrrrrrr@{}}
\toprule
Model & Run & MedQA & MedMCQA & PubMedQA & MMLU-Med & MedXpertQA & \thead{Mean\\gain} \\
\midrule
\multirow{3}{*}{MedGemma-27B-it} & Q4\_K\_M, GB10 & 69.8 & 62.1 & 76.1 & 79.2 & 16.6 & \\
 & bf16, H200, direct & 66.3 & 60.7 & 77.2 & 81.0 & 12.3 & -1.3 \\
 & bf16, H200, CoT & 83.1 & 68.5 & 76.9 & 87.6 & 21.2 & +6.7 \\
\midrule
\multirow{3}{*}{MedGemma-4B-it} & Q4\_K\_M, GB10 & 51.8 & 47.5 & 71.4 & 60.6 & 12.7 & \\
 & bf16, H200, direct & 55.2 & 52.0 & 74.9 & 64.9 & 13.2 & +3.3 \\
 & bf16, H200, CoT & 66.1 & 56.1 & 75.8 & 71.3 & 14.4 & +7.9 \\
\midrule
\multirow{3}{*}{Med42-8B (v2)} & Q4\_K\_M, GB10 & 60.4 & 57.9 & 75.6 & 71.7 & 12.4 & \\
 & bf16, H200, direct & 62.1 & 60.1 & 77.2 & 72.9 & 15.1 & +1.9 \\
 & bf16, H200, CoT & 62.7 & 58.6 & 57.9 & 73.7 & 13.8 & -2.2 \\
\midrule
\multirow{3}{*}{JSL-MedLlama-3-8B-v2.0} & Q4\_K\_M, GB10 & 27.6 & 31.0 & 36.0 & 21.6 & 9.3 & \\
 & bf16, H200, direct & 45.0 & 56.5 & 77.2 & 67.9 & 10.0 & +26.2 \\
 & bf16, H200, CoT & 58.8 & 55.5 & 76.3 & 72.9 & 12.5 & +30.1 \\
\midrule
\multirow{3}{*}{MedGemma-1.5-4B-it} & Q4\_K\_M, GB10 & 29.8 & 33.4 & 58.0 & 23.5 & 9.1 & \\
 & bf16, H200, direct & 17.8 & 40.7 & 56.1 & 41.5 & 1.7 & +0.8 \\
 & bf16, H200, CoT & 46.9 & 50.5 & 60.9 & 64.7 & 10.2 & +15.9 \\
\midrule
\multirow{3}{*}{OpenBioLLM-8B} & Q4\_K\_M, GB10 & 24.5 & 25.4 & 44.7 & 23.0 & 9.1 & \\
 & bf16, H200, direct & 43.7 & 40.3 & 44.4 & 58.2 & 9.1 & +13.8 \\
 & bf16, H200, CoT & 49.9 & 48.0 & 52.6 & 67.5 & 11.8 & +20.6 \\
\midrule
\multirow{3}{*}{MedLlama2-7B} & Q4\_K\_M, GB10 & 27.0 & 27.4 & 59.8 & 25.6 & 9.2 & \\
 & bf16, H200, direct & 16.4 & 33.2 & 68.0 & 40.3 & 4.1 & +2.6 \\
 & bf16, H200, CoT & 38.2 & 34.9 & 69.9 & 48.3 & 10.5 & +10.6 \\
\bottomrule
\end{tabular}
}
\end{table}

\section{Image results}
\label{sec:image}

\subsection{All datasets}

Table~\ref{tab:img-all} gives the seven vision models on the six image sets, sorted by the OmniMedVQA first-token score. Lingshu-32B is clearly the best model on the radiology and pathology sets: SLAKE 84.9, VQA-RAD 67.4 and PathVQA 59.9, where the next best SLAKE score is 60.3. HuatuoGPT-3-27B leads OmniMedVQA (82.9) and, with reasoning, is second on MedXpertQA-MM (49.7). Gemma-4-31B-it leads the two expert multiple-choice sets once it is allowed to reason: PMC-VQA 66.2 and MedXpertQA-MM 55.5, against 2.4 on MedXpertQA-MM under the direct protocol, where it explains instead of answering. Two OmniMedVQA strict scores are formatting failures, not knowledge failures: MedGemma-1.5-4B (19.1) and Qwen3.5-4B (27.6) describe the image before answering and are cut at 32 tokens; their first-token scores are 68.2 and 74.2, and on the replies that could be read they score 86\pct{} and 78\pct{}.

\begin{table}[htbp]
\centering\footnotesize\setlength{\tabcolsep}{3pt}
\caption{Image question accuracy in \pct{}. The first five score columns use the direct protocol; the two CoT columns use the step-by-step protocol on the two multiple-choice sets. Sorted by OmniMedVQA first-token score.}
\label{tab:img-all}
\fit{\begin{tabular}{@{}lllrrrrrrrrr@{}}
\toprule
Model & Origin & Type & VQA-RAD & SLAKE & PathVQA & PMC-VQA & \thead{MedXpertQA\\MM} & \thead{PMC-VQA\\CoT} & \thead{MedXpertQA\\MM, CoT} & \thead{OmniMedVQA\\strict} & \thead{OmniMedVQA\\first-token} \\
\midrule
HuatuoGPT-3-27B & China & Medical & 54.5 & 60.3 & 38.9 & 59.3 & 33.2 & 62.7 & 49.7 & 82.9 & 82.9 \\
Lingshu-32B & China & Medical & 67.4 & 84.9 & 59.9 & 64.9 & 31.1 & 60.1 & 31.5 & 82.4 & 82.4 \\
Gemma-4-31B-it & USA & Non-medical & 51.4 & 51.5 & 36.8 & 45.1 & 2.4 & 66.2 & 55.5 & 77.0 & 81.5 \\
Qwen3.5-4B & China & Non-medical & 55.6 & 59.6 & 39.8 & 16.2 & 0.7 & 60.0 & 35.5 & 27.6 & 74.2 \\
MedGemma-4B-it & USA & Medical & 56.8 & 52.8 & 33.1 & 46.2 & 24.1 & 46.5 & 25.4 & 70.3 & 70.4 \\
MedGemma-27B-it & USA & Medical & 50.5 & 50.0 & 32.1 & 38.2 & 4.3 & 47.1 & 33.2 & 61.9 & 68.6 \\
MedGemma-1.5-4B-it & USA & Medical & 41.2 & 35.7 & 27.7 & 17.6 & 0.2 & 45.0 & 9.6 & 19.1 & 68.2 \\
\bottomrule
\end{tabular}
}
\end{table}

\subsection{Closed and open answers}

Table~\ref{tab:img-split} splits VQA-RAD, SLAKE and PathVQA into yes/no and open short-answer questions. On yes/no questions most models sit between 70 and 90\pct{}; on open answers, scored by exact match, scores fall to between 1 and 46\pct{}, and on PathVQA, whose open answers are long medical phrases, to single digits for all but Lingshu. Lingshu-32B was trained on the training splits of these three datasets, which helps its open-answer wording (SLAKE open 81.2, PathVQA open 30.0). A correct but differently worded answer scores zero under exact match, so open scores understate every model's ability equally; papers that use a language-model judge report much higher open-answer numbers (Section~\ref{sec:img-pub}).

\begin{table}[htbp]
\centering\footnotesize\setlength{\tabcolsep}{4pt}
\caption{Radiology and pathology VQA split into yes/no and open short-answer questions, in \pct{} (direct protocol, exact match for open answers).}
\label{tab:img-split}
\fit{\begin{tabular}{@{}lllrrrrrr@{}}
\toprule
Model & Origin & Type & \thead{VQA-RAD\\yes/no} & \thead{VQA-RAD\\open} & \thead{SLAKE\\yes/no} & \thead{SLAKE\\open} & \thead{PathVQA\\yes/no} & \thead{PathVQA\\open} \\
\midrule
HuatuoGPT-3-27B & China & Medical & 75.7 & 28.0 & 82.8 & 45.9 & 71.5 & 6.2 \\
Lingshu-32B & China & Medical & 88.0 & 41.5 & 91.3 & 81.2 & 89.9 & 30.0 \\
Gemma-4-31B-it & USA & Non-medical & 75.3 & 21.5 & 87.0 & 29.8 & 70.3 & 3.2 \\
Qwen3.5-4B & China & Non-medical & 77.7 & 28.0 & 85.6 & 44.0 & 75.3 & 4.3 \\
MedGemma-4B-it & USA & Medical & 80.9 & 26.5 & 86.8 & 32.1 & 64.4 & 1.7 \\
MedGemma-27B-it & USA & Medical & 76.9 & 17.5 & 82.2 & 29.8 & 63.1 & 1.1 \\
MedGemma-1.5-4B-it & USA & Medical & 62.1 & 15.0 & 77.8 & 12.9 & 54.9 & 0.6 \\
\bottomrule
\end{tabular}
}
\end{table}

\subsection{OmniMedVQA in detail}

Table~\ref{tab:omni} breaks the 88,996 OmniMedVQA questions down by question type, using the first-token score. Modality recognition is nearly solved by every model (82 to 99\pct{}); anatomy identification separates the models most (48 to 84\pct{}), with HuatuoGPT-3 and Lingshu ahead and MedGemma-27B-it last; lesion grading is the hardest type for most models. The errors column counts the questions the engine refused: very large ISIC 2020 dermoscopy and COVIDx chest X-ray images exceed the 12,288-token context of the Qwen-family image encoder used by HuatuoGPT-3, Lingshu and Qwen3.5, and a few items carry more than four images. Excluding them would add about one point to those three models; they are kept as wrong for fairness.

\begin{table}[htbp]
\centering\footnotesize\setlength{\tabcolsep}{3pt}
\caption{OmniMedVQA (88,996 questions) in detail, in \pct{}. Question-type columns use the first-token score. Errors = questions the engine refused, counted as wrong.}
\label{tab:omni}
\fit{\begin{tabular}{@{}lllrrrrrrrrr@{}}
\toprule
Model & Origin & Type & Strict & \thead{First\\token} & \thead{Parsed\\\%} & Errors & \thead{Modality\\recognition} & \thead{Anatomy\\identification} & \thead{Disease\\diagnosis} & \thead{Lesion\\grading} & \thead{Other\\biological attributes} \\
\midrule
HuatuoGPT-3-27B & China & Medical & 82.9 & 82.9 & 98.9 & 922 & 96.2 & 84.0 & 80.3 & 70.4 & 82.7 \\
Lingshu-32B & China & Medical & 82.4 & 82.4 & 98.8 & 1,043 & 96.7 & 82.1 & 79.5 & 77.6 & 85.8 \\
Gemma-4-31B-it & USA & Non-medical & 77.0 & 81.5 & 92.2 & 0 & 99.3 & 79.5 & 78.7 & 65.5 & 86.5 \\
Qwen3.5-4B & China & Non-medical & 27.6 & 74.2 & 35.4 & 922 & 96.7 & 76.3 & 69.8 & 52.8 & 72.4 \\
MedGemma-4B-it & USA & Medical & 70.3 & 70.4 & 99.5 & 0 & 96.8 & 57.8 & 68.4 & 70.4 & 74.8 \\
MedGemma-27B-it & USA & Medical & 61.9 & 68.6 & 87.3 & 0 & 88.0 & 47.9 & 69.7 & 77.7 & 77.4 \\
MedGemma-1.5-4B-it & USA & Medical & 19.1 & 68.2 & 22.3 & 0 & 81.9 & 61.4 & 67.6 & 67.6 & 63.7 \\
\bottomrule
\end{tabular}
}
\end{table}

\subsection{Comparison with published image scores}
\label{sec:img-pub}

Table~\ref{tab:img-pub} lists every published image score we could find for these models and whether the comparison is like for like. Where it is, the harness agrees: Lingshu-32B scores 82.4 against a published 83.4 on OmniMedVQA and 31.5 against 30.9 on MedXpertQA-MM; MedGemma-4B-it scores 70.4 against 69.8 on OmniMedVQA and 86.8 against 87.6 on SLAKE closed questions. Our VQA-RAD closed scores are about 10 points above the Google cards for both MedGemma models. Where the published number uses a different rule the gap is not a quality gap: Lingshu's paper has GPT-4.1 judge the open answers of SLAKE, VQA-RAD and PathVQA, and the Gemma 4 and Qwen cards use their thinking mode with long budgets, which we switched off. No published image scores exist for HuatuoGPT-3-27B, and Google reports neither PathVQA nor OmniMedVQA for any MedGemma model.

\begin{table}[htbp]
\centering\footnotesize\setlength{\tabcolsep}{4pt}
\caption{This run against the published image scores, in \pct{}. Diff. = this run minus published, in points. Rows marked No are listed for completeness: the published number uses a different scoring rule.}
\label{tab:img-pub}
\fit{\begin{tabular}{@{}llrlrlrl@{}}
\toprule
Model & Benchmark & This run & Our protocol & Published & Source & Diff. & Like for like? \\
\midrule
Lingshu-32B & OmniMedVQA & 82.4 & first-token & 83.4 & Lingshu paper & -1.0 & Yes \\
Lingshu-32B & PMC-VQA & 64.9 & direct & 57.9 & Lingshu paper & +7.0 & Yes \\
Lingshu-32B & MedXpertQA-MM & 31.5 & CoT & 30.9 & Lingshu paper & +0.6 & Yes \\
Lingshu-32B & SLAKE & 84.9 & exact match & 89.2 & Lingshu paper & -4.3 & No: GPT-4.1 judges open answers \\
Lingshu-32B & VQA-RAD & 67.4 & exact match & 76.5 & Lingshu paper & -9.1 & No: GPT-4.1 judges open answers \\
Lingshu-32B & PathVQA & 59.9 & exact match & 65.9 & Lingshu paper & -6.0 & No: GPT-4.1 judges open answers \\
MedGemma-4B-it & OmniMedVQA & 70.4 & first-token & 69.8 & Lingshu paper & +0.6 & Yes \\
MedGemma-4B-it & PMC-VQA & 46.2 & direct & 49.9 & Lingshu paper & -3.7 & Yes \\
MedGemma-4B-it & SLAKE closed (yes/no) & 86.8 & yes/no & 87.6 & Google card & -0.8 & Yes \\
MedGemma-4B-it & VQA-RAD closed (yes/no) & 80.9 & yes/no & 69.1 & Google card & +11.8 & Yes \\
MedGemma-4B-it & MedXpertQA-MM & 25.4 & CoT & 18.8 & Google card & +6.6 & Yes \\
MedGemma-27B-it & SLAKE closed (yes/no) & 82.2 & yes/no & 85.9 & Google 1.5 card & -3.7 & Yes \\
MedGemma-27B-it & VQA-RAD closed (yes/no) & 76.9 & yes/no & 67.1 & Google 1.5 card & +9.8 & Yes \\
MedGemma-27B-it & MedXpertQA-MM & 33.2 & CoT & 26.8 & Google card & +6.4 & Roughly: card mixes text and image items \\
MedGemma-1.5-4B-it & SLAKE closed (yes/no) & 77.8 & yes/no & 82.8 & Google 1.5 card & -5.0 & Yes \\
MedGemma-1.5-4B-it & VQA-RAD closed (yes/no) & 62.1 & yes/no & 70.2 & Google 1.5 card & -8.1 & Yes \\
MedGemma-1.5-4B-it & MedXpertQA-MM & 9.6 & CoT & 20.9 & Google 1.5 card & -11.3 & Yes \\
Gemma-4-31B-it & MedXpertQA-MM & 55.5 & CoT & 61.3 & Gemma 4 report & -5.8 & Roughly: report uses thinking mode \\
Qwen3.5-4B & PMC-VQA & 60.0 & CoT & 55.5 & Qwen card & +4.5 & Roughly: card uses thinking mode \\
Qwen3.5-4B & MedXpertQA-MM & 35.5 & CoT & 42.9 & Qwen card & -7.4 & Roughly: card uses thinking mode \\
Qwen3.5-4B & SLAKE & 59.6 & exact match & 76.1 & Qwen card & -16.5 & No: scoring not stated \\
\bottomrule
\end{tabular}
}
\end{table}

\section{Throughput}
\label{sec:throughput}

Table~\ref{tab:throughput} gives the load time, the run time over the five text sets and the mean number of generated tokens per question for each model and protocol, on the power-capped H200 with vLLM batching. Under the direct protocol every model finishes the 9,995 questions in 2 to 11 minutes. The step-by-step protocol costs 5 to 25 times more, in proportion to how much each model writes: MedGemma-27B-text-it writes 1,342 tokens per question on average and takes 2.6 hours, Lingshu-32B writes 271 and takes 31 minutes. These timings are for batched evaluation, not for the interactive latency of the deployed system (Section~\ref{sec:speed}).

\begin{table}[htbp]
\centering\small\setlength{\tabcolsep}{6pt}
\caption{Load and run time over the five text sets (9,995 questions) and mean generated tokens per question, on one H200 at 300\,W.}
\label{tab:throughput}
\begin{tabular}{@{}llrrr@{}}
\toprule
Model & Protocol & Load (s) & Run (s) & Mean generated tokens \\
\midrule
Gemma-4-31B-it & direct & 182 & 677 & 9 \\
Gemma-4-31B-it & step-by-step & 116 & 4749 & 423 \\
HuatuoGPT-3-27B & direct & 128 & 528 & 2 \\
HuatuoGPT-3-27B & step-by-step & 83 & 5873 & 872 \\
Qwen3.5-4B & direct & 78 & 149 & 2 \\
Qwen3.5-4B & step-by-step & 57 & 1263 & 805 \\
MedGemma-27B-text-it & direct & 118 & 385 & 20 \\
MedGemma-27B-text-it & step-by-step & 68 & 9225 & 1342 \\
MedGemma-27B-it & direct & 116 & 351 & 5 \\
MedGemma-27B-it & step-by-step & 59 & 2415 & 1147 \\
Baichuan-M2-32B & direct & 86 & 542 & 6 \\
Baichuan-M2-32B & step-by-step & 55 & 3166 & 482 \\
Lingshu-32B & direct & 91 & 548 & 3 \\
Lingshu-32B & step-by-step & 60 & 1833 & 271 \\
MedGemma-4B-it & direct & 84 & 114 & 3 \\
MedGemma-4B-it & step-by-step & 33 & 576 & 260 \\
Med42-8B (v2) & direct & 39 & 98 & 6 \\
Med42-8B (v2) & step-by-step & 38 & 386 & 236 \\
JSL-MedLlama-3-8B-v2.0 & direct & 55 & 175 & 32 \\
JSL-MedLlama-3-8B-v2.0 & step-by-step & 38 & 596 & 206 \\
MedGemma-1.5-4B-it & direct & 74 & 148 & 13 \\
MedGemma-1.5-4B-it & step-by-step & 33 & 1882 & 1278 \\
OpenBioLLM-8B & direct & 57 & 165 & 27 \\
OpenBioLLM-8B & step-by-step & 40 & 395 & 31 \\
MedLlama2-7B & direct & 43 & 162 & 27 \\
MedLlama2-7B & step-by-step & 34 & 605 & 149 \\
\bottomrule
\end{tabular}

\end{table}

\section{Grounded evaluations for the Tumour Board}
\label{sec:grounded}

Examinations measure knowledge. Vaidyx-TB needs something else: a reader that transcribes a scanned page without inventing a number, an extractor that fills the Board's fields only with values it can quote, a reasoner that answers a question from the Case File or says that the file does not contain the answer, and a pipeline that tells a draft report from a final one. These are measured on a synthetic corpus of Case Files with an answer key (no real patient data) and, where real text is needed, on public de-identified pathology reports from The Cancer Genome Atlas (TCGA). All of these evaluations run through the same grounding layer as the product: a value or sentence is kept only when its verbatim quote is found in the source document; otherwise it is dropped or marked for review.

\subsection{Reader and extractor bake-off}

The corpus has three Case Files and 31 documents (PDF text layers, five scanned pages and one text file) with an answer key. The \emph{reader} transcribes the five scans and is scored by word recall against the reference text and by the number of invented numbers; the \emph{extractor} classifies the document type and fills 49 fields and is scored on types, fields and fabrications (values with no quote). Table~\ref{tab:bakeoff} shows the result. Qwen3.5-4B is the most accurate and by far the fastest: perfect recall, no invented number, every type and field, 1.6\,s for the five scans and 4.8\,s for the 31 documents, at 9\,GB of weights. It is Vaidyx-TB's reader and extractor, labelled Chinese and non-medical throughout the product. MedGemma-27B-it invented one number in transcription.

\begin{table}[htbp]
\centering\footnotesize\setlength{\tabcolsep}{4pt}
\caption{Reader and extractor bake-off on the synthetic corpus (3 Case Files, 31 documents, 5 scans). Fabricated = extracted values with no verbatim quote. Times on the H200 at 300\,W.}
\label{tab:bakeoff}
\fit{\begin{tabular}{@{}lllrrrrrrr@{}}
\toprule
Model & Origin & Type & \thead{Reader\\word recall} & \thead{Invented\\numbers} & \thead{Reader s\\(5 scans)} & \thead{Doc.\\types} & Fields & Fabricated & \thead{Extractor s\\(31 docs)} \\
\midrule
Qwen3.5-4B & China & Non-medical & 1.000 & 0 & 1.6 & 31/31 & 49/49 & 0 & 4.8 \\
Gemma-4-31B-it & USA & Non-medical & 1.000 & 0 & 11.1 & 30/31 & 49/49 & 0 & 28.1 \\
MedGemma-27B-it & USA & Medical & 0.976 & 1 & 7.9 & 31/31 & 49/49 & 0 & 23.0 \\
Lingshu-32B & China & Medical & 0.994 & 0 & 11.3 & 31/31 & 48/49 & 0 & 22.5 \\
HuatuoGPT-3-27B & China & Medical & 1.000 & 0 & 8.4 & 30/31 & 49/49 & 0 & 27.2 \\
MedGemma-27B-text-it & USA & Medical & text only & -- & -- & 30/31 & 49/49 & 0 & 23.3 \\
\bottomrule
\end{tabular}
}
\end{table}

\subsection{Answerability of questions to the Case File}

The reasoner answers questions about a prepared Case File. Thirty-three probes cover answerable questions, absent values (fabrication traps, where the right answer is that the file does not say) and out-of-scope requests (a treatment recommendation, which must be declined). Table~\ref{tab:answerability} gives the three candidates. Gemma-4-31B-it passes all 33 and is the fastest, so it is the reasoner, labelled non-medical everywhere in the product; HuatuoGPT-3-27B and MedGemma-27B-it are one configuration switch away. All three produced zero fabrications, which is the grounding layer's doing rather than the models'. The same 33 probes are re-run at every release on freshly prepared cases and still pass with zero fabrications.

\begin{table}[htbp]
\centering\small
\caption{Ask-the-Case-File answerability, 33 probes. Latency is the mean per answer on the H200 at 300\,W; Gemma's figure is after the parallel scope check.}
\label{tab:answerability}
\begin{tabular}{@{}lllrrr@{}}
\toprule
Reasoner & Origin & Type & Pass & Fabrications & Mean latency (s) \\
\midrule
Gemma-4-31B-it & USA & Non-medical & 33/33 & 0 & 3.0 \\
HuatuoGPT-3-27B & China & Medical & 32/33 & 0 & 8.0 \\
MedGemma-27B-it & USA & Medical & 31/33 & 0 & 8.4 \\
\bottomrule
\end{tabular}

\end{table}

\subsection{Real pathology reports}

To check the pipeline on real text, 30 pathology reports from TCGA (OCR text of the public de-identified PDFs) were processed with the shipped reader and extractor and compared with the TCGA clinical ground truth. The extractor produced 166 values; none was ungrounded. Agreement with the ground truth was ER 7/8, PR 7/8, histology 27/29 and pT as recorded 24/26 (Table~\ref{tab:tcga}). Every disagreement is a metric artefact rather than a wrong reading: the report says ``95\pct{} nuclear staining with moderate intensity'' where the ground truth says ER positive, the report states a stage group where the ground truth gives a pT, and the ground truth says ``Other, specify''. The run took 6.7\,s for the 30 reports.

\begin{table}[htbp]
\centering\small
\caption{Thirty real TCGA pathology reports through the shipped pipeline: extracted values, grounding and agreement with the clinical ground truth.}
\label{tab:tcga}
\begin{tabular}{@{}rrrrrrr@{}}
\toprule
Reports & Extracted values & Ungrounded & ER & PR & Histology & pT (as recorded) \\
\midrule
30 & 166 & 0 & 7/8 & 7/8 & 27/29 & 24/26 \\
\bottomrule
\end{tabular}

\end{table}

\subsection{Report status on 9,523 real reports}

Whether a report is final decides whether a pre-authorisation requirement is met. The first status detector was tuned on the synthetic corpus, whose reports carry a ``Report status:'' line. Measured on 9,523 real TCGA pathology reports \citep{kefeli2024tcga} (Table~\ref{tab:status}), 26\pct{} say ``FINAL DIAGNOSIS'', 15\pct{} ``ADDENDUM'', 14\pct{} ``verified'' or ``reviewed by'', 5\pct{} ``signed'', 3\pct{} carry an end-of-report marker, and 61\pct{} carry no explicit status or signature evidence at all, because de-identification removes signatures. A hard rule would have stalled most real Case Files. The shipped policy follows the FHIR DiagnosticReport status vocabulary: an explicit negative (draft, preliminary, provisional, interim, cancelled, signature pending, unsigned) never counts; an explicit positive (a final header, an end-of-report marker, an electronic signature, verified or approved by, a registration number) counts as final, as do addendum, amended and corrected; no marker at all is ``unstated'', which counts when a signature line is present and otherwise turns the requirement into ``confirm final version'': neither blocked nor passed, and a clinician can attest with one audited click. ``Working Draft'' appears inside finalised TCGA reports as a template artefact, so the word draft alone only adds a note.

\begin{table}[htbp]
\centering\small
\caption{Status evidence found in 9,523 real TCGA pathology reports (a report can carry more than one marker).}
\label{tab:status}
\begin{tabular}{@{}lr@{}}
\toprule
Evidence in the report & Share of the 9,523 reports \\
\midrule
FINAL DIAGNOSIS & 26 \% \\
ADDENDUM & 15 \% \\
verified / reviewed by & 14 \% \\
signed & 5 \% \\
end of report & 3 \% \\
no explicit status or signature & 61 \% \\
\bottomrule
\end{tabular}

\end{table}

\subsection{Summary omission}

After full processing of a synthetic Case File, the Board summary and the verified primary-fields table are checked against the answer key (a recall metric in the style of TBFact): 47 of 48 answer-key fields appear in the verified fields or in a cited summary sentence. The one miss is a negative finding in a discharge summary (``IV antibiotics were not required'') that the Board summary did not restate. The check is repeated at every release.

\subsection{Scanned pages: two reads and an independent OCR engine}
\label{sec:ocr}

Generative readers silently ``correct'' digits. Vaidyx-TB therefore reads every scanned page twice (at full resolution and at a 70\pct{} downscale), aligns the two reads line by line and marks any number that the reads disagree on; a value quoted from such a line needs review. Both reads come from the same model, so a digit it corrects the same way twice would pass. An independent classical OCR engine was added as a third read of the numbers only. Table~\ref{tab:ocr} compares the two candidate engines on the 111 synthetic scanned pages, which have an exact reference text, on CPU only. Tesseract 5.5.3 \citep{smith2007tesseract} with page segmentation mode 4 reaches 98.6\pct{} number recall and 98.4\pct{} precision at 0.395\,s a page on one thread; docTR 1.1.0 reaches 97.7\pct{} and 95.8\pct{} once its words are regrouped into rows (17.8\pct{} on its native block lines, which split every table cell) at 1.276\,s a page on eight threads. Tesseract was integrated: it runs in its own environment with no Python package in the application, one thread per page, in parallel with the two model reads, and it never changes the transcription.

The check itself was tuned to be quiet. Table~\ref{tab:ocr-check} shows its false alarms with the reference text as the transcription and what it catches of planted one-digit errors (one per line with a compared number: a substituted digit, an inserted digit or a deleted digit). The two-read rule applied as is raises 54 false alarms on these pages (lost decimal points, cut leading digits, list markers, letters read as digits). Comparing digits only, requiring the engine to read another number with the same count of digits on the partner line, and keeping a flag only when a second engine read at 1.5$\times$ scale doubts the same number brings the false alarms to zero while still catching 834 of 839 substitutions (99.4\pct{}); the shipped rule adds insertions and deletions and catches 90.3\pct{} and 87.9\pct{} of those. A flagged line puts every value quoted from it into review. Limits: the corpus is synthetic (one scan rasterised at about 86\,dpi with noise, tilt and JPEG compression, one phone photograph); real scans with faint stamps and handwriting will give Tesseract more misses, and a missed number never raises a flag, so the check gets quieter, not noisier.

\begin{table}[htbp]
\centering\small
\caption{Independent OCR engines on 111 synthetic scanned pages with an exact reference text, CPU only. Recall = reference numbers found on the aligned engine line; precision = engine numbers found on the aligned reference line.}
\label{tab:ocr}
\fit{\begin{tabular}{@{}lrrrr@{}}
\toprule
Engine & Number recall \% & Number precision \% & s / page (median) & Threads per page \\
\midrule
Tesseract 5.5.3 (psm 4) & 98.6 & 98.4 & 0.395 & 1 \\
docTR 1.1.0 (rows) & 97.7 & 95.8 & 1.276 & 8 \\
docTR 1.1.0 (native block lines) & 17.8 & 17.6 & 1.276 & 8 \\
\bottomrule
\end{tabular}
}
\end{table}

\begin{table}[htbp]
\centering\small
\caption{The OCR number check with Tesseract: false alarms with the reference text as the transcription (111 pages) and planted one-digit errors caught (839 substitutions, 839 insertions, 735 deletions).}
\label{tab:ocr-check}
\fit{\begin{tabular}{@{}lrrrr@{}}
\toprule
Check & False alarms & Substitutions caught \% & Insertions caught \% & Deletions caught \% \\
\midrule
3.3 rule as is & 54 & 99.6 & 99.8 & 99.7 \\
Substitution, one read & 2 & 99.4 & 2.0 & 2.7 \\
Substitution, confirming re-read & 0 & 99.4 & 0.2 & 0.3 \\
Shipped: substitution, insertion, deletion, confirming re-read & 0 & 99.4 & 90.3 & 87.9 \\
\bottomrule
\end{tabular}
}
\end{table}

\subsection{Interactive speed}
\label{sec:speed}

On the power-capped H200, Gemma-4-31B-it decodes at about 30 tokens a second. Prompt-lookup (n-gram) speculative decoding on the reasoner gives byte-identical greedy output, which was verified, and about three times faster generation: a specialty brief falls from 37.6\,s to 12.4\,s and a Case File summary from 63.8\,s to 20.8\,s. Ingestion of 27 documents including five scans takes about 6\,s with all documents in parallel; the instant value search takes about 2\,ms and a guideline query about 120\,ms over a library of 97 files and 5,759 passages, India first (NCG, ICMR, PM-JAY, ICRI/IRIA), then ESMO, CAP/ASCO-CAP, ACR, RECIST, NCI PDQ, RCPath, RCR and NICE; passages are shown as reference text, never as a recommendation.

\section{Discussion}

\paragraph{The harness is like for like.} Where the published protocol matches ours (generated answers, no thinking mode, the same scoring rule), our scores agree with the cards within a few points in both directions (Tables~\ref{tab:text-pub} and \ref{tab:img-pub}). The disagreements are explained by protocol: answer likelihood against generated answers, a language-model judge against exact match, thinking mode against a visible chain of thought, and one model that runs out of its 2,048-token budget. The numbers in this report can therefore be read as like for like with each other, which is what a model choice needs.

\paragraph{Model generation matters more than size or medical fine-tuning.} The two best text models and the best small model are general models or recent medical models built on current base weights. Qwen3.5-4B, a 4B general model, matches a 27B medical model on text and beats every 7B to 8B medical fine-tune of the previous generation by 10 to 28 points. The older fine-tunes (MedLlama2, OpenBioLLM, JSL-MedLlama, Med42) are superseded and need not be benchmarked again.

\paragraph{The hard sets separate the models.} On MedQA and MMLU-Med the top models are close to the ceiling. MedXpertQA (10 options, filtered to be hard) and MedMCQA (the Indian entrance examinations) spread them out: 55.1 against at most 40.4 on MedXpertQA, and 9 to 16 points lower on MedMCQA than on MedQA for the seven strongest models. An India-specific test is the obvious next addition.

\paragraph{Format discipline is a separate property.} Several strong models cannot reply with a letter when asked to (MedGemma-27B-text-it, MedGemma-1.5-4B, Qwen3.5-4B and MedGemma-1.5-4B on images). The first-token and parse-rate columns separate this from knowledge. For a system that asks structured questions, the direct-answer columns are as informative as the step-by-step ones, and they favour HuatuoGPT-3-27B and Lingshu-32B.

\paragraph{Benchmarks do not measure grounding.} All three candidate reasoners produced zero fabrications on the 33 probes, and the 166 values from real pathology reports were all grounded, because the grounding layer, not the model, decides what is shown. The benchmark picks the model with the most knowledge and the best reasoning; the grounded evaluations decide whether the product may show its output. Both are needed, and the second is the one that keeps the system within its declared scope.

\section{Caveats}

\begin{itemize}
  \item \textbf{Thinking mode was off.} HuatuoGPT-3, Baichuan-M2 and Qwen3.5 publish results with thinking mode on and long budgets (up to 32k tokens). Here every model got the same visible step-by-step prompt and 2,048 tokens, so these three may score higher in their own settings.
  \item \textbf{Open image answers use exact match.} Correct but differently worded short answers, on PathVQA especially, score zero. Papers often use a language-model judge (Lingshu uses GPT-4.1), which gives much higher open-answer numbers.
  \item \textbf{Engine errors count as wrong.} HuatuoGPT-3 (922), Lingshu (1,043) and Qwen3.5 (922) had OmniMedVQA questions refused because large images exceed the context window. Excluding them adds about one point.
  \item \textbf{The short-answer format penalises explaining models.} Use the step-by-step columns to judge ability and the direct columns to judge format discipline.
  \item \textbf{Baichuan-M2's direct scores were re-scored} from the saved replies after a stray thinking tag broke the parser; no answer was regenerated. The parser will strip thinking tags in the next round.
  \item \textbf{Possible training overlap.} Lingshu was trained on the training splits of VQA-RAD, SLAKE and PathVQA, and newer models may have seen public test items. Benchmarks measure examination knowledge, not safety on real patient documents.
  \item \textbf{The grounded evaluations are small and mostly synthetic.} Three Case Files, 31 documents, 33 probes, 111 synthetic scanned pages, and 30 real reports for the pipeline check. They are release gates for one product, not a clinical validation, and no clinical claim is made from them. Indic-language OCR is installed but not measured.
  \item \textbf{One GPU, power-capped.} Timings are specific to an H200 at 300\,W; scores are not affected.
\end{itemize}

\section{What Vaidyx-TB uses, and the next round}

Gemma-4-31B-it (non-medical) is the reasoner: the strongest text model, with the existing verbatim-quote grounding because it is not medically validated. Qwen3.5-4B (Chinese, non-medical) is the reader and extractor: 62.4\pct{} overall at 4B, ahead of every 7B to 8B medical model, and the winner of the bake-off; MedGemma-4B-it is the best small option from outside China (50.8\pct{}). HuatuoGPT-3-27B (Chinese, medical) is the best medical model on text and OmniMedVQA and is worth a Vaidyx-specific evaluation as a second opinion where a Chinese-origin model is acceptable for the deployment. Lingshu-32B (Chinese, medical) leads radiology and pathology VQA and Gemma-4-31B-it the expert multi-image questions; neither reads a scan inside Vaidyx-TB, which only reads report documents and leaves images to the clinicians. The next benchmark round will strip thinking tags in the parser, re-run the direct image tests with 128 tokens so that explaining models are scored fairly, add an Indian-language and Indian-guideline test set, and score open image answers with a language-model judge alongside exact match.

\section{Reproducibility and sources}
\label{sec:sources}

Every model ran from its published bf16 Hugging Face checkpoint with vLLM 0.30, greedy decoding, seed 0, its own chat template and thinking mode off; datasets are the full public test splits named in Table~\ref{tab:datasets}. Every reply is stored per question with the parsed answer, so any score in this report can be recomputed from the saved outputs, and the Baichuan-M2 re-scoring was done that way. Published numbers were opened on 9 October 2026 from: the MedGemma model card (27B, 27B-text and 4B tables) and the MedGemma 1.5 4B card (closed-question accuracy for 4B and 27B) \citep{sellergren2025medgemma}; the Lingshu-32B card and technical report, Tables I and II \citep{xu2025lingshu}; the Gemma 4 31B card and technical report, Tables 5, 6 and 12 \citep{gemma4}; the Qwen3.5-4B card \citep{qwen35}; the Llama3-Med42-8B card \citep{christophe2024med42}; the Llama3-OpenBioLLM-8B card \citep{openbiollm}; the JSL-MedLlama-3-8B-v2.0 card \citep{jslmedllama}; the HuatuoGPT-3 repository \citep{huatuogpt3}; and the Baichuan-M2-32B card \citep{baichuanm2}. The earlier 4-bit figures are from the Vaidyx benchmark report of September 2026 (Q4\_K\_M, GB10, Ollama). The benchmarks are MedQA \citep{jin2021medqa}, MedMCQA \citep{pal2022medmcqa}, PubMedQA \citep{jin2019pubmedqa}, MMLU \citep{hendrycks2021mmlu}, MedXpertQA \citep{zuo2025medxpertqa}, VQA-RAD \citep{lau2018vqarad}, SLAKE \citep{liu2021slake}, PathVQA \citep{he2020pathvqa}, PMC-VQA \citep{zhang2023pmcvqa} and OmniMedVQA \citep{hu2024omnimedvqa}. The real reports are the TCGA pathology reports \citep{kefeli2024tcga}. The OCR engines are Tesseract 5.5.3 \citep{smith2007tesseract} and docTR 1.1.0 \citep{doctr}.

\bibliographystyle{unsrtnat}
\bibliography{refs}

\end{document}
